\chapter{The next C2 boundary: retain the fourth event current}
\label{component:p53752031d16c}
\noindent\textbf{Component }\path{p53752031d16c}\par
{\small\noindent Source: \path{CORE/prior/inherited/c2_y_current.tex}\par}
{\small\noindent SHA256: \path{53752031d16c6691f061be1e7f5980fbd300cd3e60d5c03e599ed99f4179408a}\par}
{\small\noindent Inventory use: Fourth centered current and bounded retained remainder, exact six-root zero relation/correction, full W-unit state identity. Its N-unit theorem is not imported on the radial domain.\par}
\medskip
\section{The next C2 boundary: retain the fourth event current}
\label{p53752031d16c:sec:c2-y-current}

The gap $h=b-a^2t$ need not stay a unit.  When it becomes
infinitesimal, the fourth event is strong.  Its centered current supplies
an exact zero relation which repairs the preceding Gram matrix.

\subsection{Source positivity and the precise remaining corner}
Keep the exact source \eqref{p9f1235511111:eq:c2u-source}, with $c>0$ infinitesimal,
$r,g,j>0$ bounded, and $a,1-r,t$ positive bounded units.  This section
assumes
\begin{equation}
h=b-a^2t>0\text{ infinitesimal},\qquad N\text{ a positive bounded unit}.
\label{p53752031d16c:eq:c2y-domain}
\end{equation}
There is a useful exact positive identity:
\begin{equation}
N=(1+c^2)(b+r)h+a^2\bigl[g+j+rt(1+a^2c^2)\bigr].
\label{p53752031d16c:eq:c2y-N-positive}
\end{equation}
It follows by writing $b+a^2(1-t)=h+a^2$ in the definition of $N$
and using the definitions of $b,t$.  Thus all source coefficients are
positive.  Under the other unit hypotheses, $N$ is a unit exactly when
at least one of $h,r,g,j$ has positive residue.  In
\eqref{p53752031d16c:eq:c2y-domain}, therefore, at least one of $r,g,j$ is ordinary.
As before, $re_1,ge_2,je_5,e_3$ are positive bounded units, and now
\begin{equation}
U_4=he_4=\frac{a^2A_c(b+r)(A_c+a^2c^2t)}{2\Delta},\qquad
\frac{\alpha}{c^4}=\frac{(b+r)N}{2(1-c^2h)\Delta}
\label{p53752031d16c:eq:c2y-units}
\end{equation}
are positive bounded units.  Hence $e_4$ is strong and
$1/\sqrt{e_4}$ is infinitesimal.

\begin{theorem}[Full parent at the fourth-event boundary]
\label{p53752031d16c:thm:c2-y-current}
For the exact source under \eqref{p53752031d16c:eq:c2y-domain}, every independent
testing pair $u,v$ that is infinitesimal or zero, and every real or
complex projective state, the entire parent \eqref{p9f1235511111:eq:c2u-parent} is a
finite sum of squares in $K_c\otimes_{\mathbb R}K_t$.
\end{theorem}

\subsection{The exact zero relation}
First take $Z=1$ and bounded $W\geq0$.  In addition to the three currents
already defined, put
\[
H_4=Y-\frac ha u.
\]
This current vanishes at the exact stationary contact, since
$W_c-at=h/a$.  Retain the extra packet $4e_4D_uD_vH_4^2$.
Its difference from the original fourth event is
\[
4e_4D_uD_v(Y^2-H_4^2)
=\frac{8U_4}{a}D_uD_vuH_4+
\frac{4hU_4}{a^2}D_uD_vu^2.
\]
Thus the entire remainder after retaining four packets has bounded
coefficients.  Its exact contact, parity, quadratic state degree, and
value $4\alpha$ at $u=v=H_1=H_5=0$ are unchanged.  The complete
monic-division argument in Section~\ref{p9f1235511111:sec:c2-unified} therefore gives
a bounded quadratic form on $(H_1,H_5,R,S,T)$.  Eliminate $S$ by
\eqref{p9f1235511111:eq:c2u-eliminate-S} and restore the four retained packets.  This
gives a bounded exact Gram on the six roots
\begin{equation}
K_1=\sqrt{e_1}H_1,\quad K_2=\sqrt{e_2}H_2,\quad
K_4=\sqrt{e_4}H_4,\quad K_5=\sqrt{e_5}H_5,\quad R,\quad T.
\label{p53752031d16c:eq:c2y-roots}
\end{equation}

There is one exact relation among these roots.  With
$K=1-j+ru^2$ and $Q$ as in Section~\ref{p9f1235511111:sec:c2-unified}, direct
substitution gives
\[
H_4=S+KR-ar uT+uvH_1-H_5.
\]
Combining this with the current formula for $S$ yields
\begin{align}
\zeta={}&aD_uH_4-H_2-(1-u^2v^2+aD_uuv)H_1
 +(aD_u-uv)H_5\notag\\
&-[D_uW+(a-uv)K]R+u[Q(u^2)+a^2D_ur]T=0.
\label{p53752031d16c:eq:c2y-zero-relation}
\end{align}
In the scaled roots every coefficient here is bounded.  Its residue is
\begin{equation}
\overline\zeta=-\frac{K_2}{\sqrt m}-\frac{K_1}{\sqrt n}
+\frac{aK_5}{\sqrt D}-(W+at_0)R,
\label{p53752031d16c:eq:c2y-zero-residue}
\end{equation}
where, here and in residue calculations below, coefficient expressions
mean their residues; $n,m,C,A,D,t_0,w_0$ are the reference quantities
defined in the preceding section.  The $K_4$ and $T$ coefficients
vanish at residue.  The coefficient of $R$ is a nonzero real number,
since $W\geq0$ and $a,t_0$ are positive units.

\subsection{The corrected residue is visibly strict}
Let $h_0=b_0-a^2t_0$.  The definitions give
\begin{equation}
h_0=(1+c^2)h+a^2c^2r>0,
\label{p53752031d16c:eq:c2y-h0}
\end{equation}
so $h_0$ is infinitesimal too.  The reference fourth weight satisfies
$h_0A>0$ with unit residue, and $U_4-h_0A$ is infinitesimal.
The ratio $e_4/A$ need not have residue one.  It is unnecessary:
the bounded retained remainder depends on $U_4,h$, while the fourth
scaled packet has coefficient $4D_uD_v$ and
$1/\sqrt{e_4}=\sqrt h/\sqrt{U_4}$ has zero residue.  These facts
justify the reference comparison even when $h\ll c^2r$.

For a transparent calculation, start from
\eqref{p9f1235511111:eq:c2u-pre-scaled}.  The linearized fourth current is
\[
H_{4,0}=S+t_0R-H_5
=\frac{(W+at_0)R+H_1+H_2-aH_5}{a}.
\]
The last term of that reference quadratic form is exactly
\[
4A\left(H_{4,0}+\frac{h_0}{a^2}R\right)^2.
\]
Subtract the fourth retained square $4AH_{4,0}^2$ and restore
$4K_4^2$.  The resulting four-hard-root part is
\[
\begin{split}
&(4+2/n)K_1^2+4K_2^2+4K_4^2
 +\bigl(4+(4C+2)/D\bigr)K_5^2\\
&\hspace{12mm}+\frac{8h_0A}{a^2}H_{4,0}R
 +\frac{4h_0^2A}{a^4}R^2.
\end{split}
\]
At residue, the last square coefficient vanishes,
$H_{4,0}=-\overline\zeta/a$, and $h_0A$ has the same residue as $U_4$.
The radial calculation is unchanged: its pivot is
$4\overline{\alpha/c^4}>0$ and all radial cross coefficients vanish.
Thus the entire Gram residue is the residue of
\begin{equation}
\mathcal D+4(\alpha/c^4)T^2-\frac{8U_4}{a^3}\zeta R,
\quad
\mathcal D=(4+2/n)K_1^2+4K_2^2+4K_4^2
+\bigl(4+(4C+2)/D\bigr)K_5^2.
\label{p53752031d16c:eq:c2y-uncorrected}
\end{equation}
All these residue comparisons use bounded normalized coefficients and
monic division, exactly as in the preceding proof.

Now add the following identically zero polynomial to the exact Gram:
\begin{equation}
\frac{8U_4}{a^3}\zeta R+\zeta^2=0.
\label{p53752031d16c:eq:c2y-zero-correction}
\end{equation}
Its coefficients are bounded, and the equality follows from the exact
relation \eqref{p53752031d16c:eq:c2y-zero-relation}.  The corrected residue is
\[
\overline{\mathcal D}
+4\overline{\alpha/c^4}\,T^2+\overline\zeta^{\,2}\succ0.
\]
Indeed the first four diagonal coefficients are at least $4$, the radial
pivot is positive, and after those five roots vanish
$\overline\zeta^{\,2}$ has positive coefficient on $R$ by
\eqref{p53752031d16c:eq:c2y-zero-residue}.  The strict-residue Gram lemma proves the
full bounded $Z=1,W\geq0$ chart.

\subsection{The other projective chart}
For $W=1$, $Z$ infinitesimal, retain the original packets in
$X,R_2,Y,R_5$ and the $\alpha$ packet.  The exact pure-testing identity
\begin{equation}
u=\frac{D_uY+[Z(1-u^2v^2)-D_uuv]X+(D_u+uvZ)R_5+ZR_2}
{D_u-D_vZ^2}
\label{p53752031d16c:eq:c2y-Wunit}
\end{equation}
follows from the two amplitude definitions.  Its denominator is a
positive pure-testing unit.  Together with \eqref{p9f1235511111:eq:c2u-wunit-root},
it expresses $u,v$ linearly in $X,R_2,Y,R_5$, with bounded testing
coefficients.  The signed-core completion in the preceding section
then makes the whole remaining part a bounded quadratic form on the
four scaled packet roots.  Its residue is
\[
(4+2/e_1)K_1^2+4K_2^2+4K_4^2+
\bigl(4+(4e_3+2)/e_5\bigr)K_5^2\succ0.
\]
This proves the chart for arbitrary $X,Y$.  The same testing-only sign
gauges, signed-core Schur complement, homogeneous restoration, and
real/imaginary splitting complete all projective states.

The only remaining corner of this particular bounded-shape source
family has $h,r,g,j$ all infinitesimal.  Then $a^2$ has residue $1$,
$N$ is infinitesimal, and the radial pivot used here vanishes.
Neither this theorem nor the preceding unit-gap theorem asserts a
certificate for that further degeneration.
